
\section{第 6 章结果的陈述}
\label{sec:6-1}

本章的目标是从滤过微分
$\Lambda_{0,\mathrm{nov}}$-模
构造谱序列，
从而证明定理 D、E、G，
继而证明定理 H、I、J、K、L。
第二作者在
\textbf{[Oh96I,II]}
中曾为单调拉格朗日子流形
构造类似的谱序列；
在该情形，本章的谱序列
与 \textbf{[Oh96I]} 的构造一致\liftIndex{I370}。
不过，两种构造有两点差别：
其一，
\textbf{[Oh96I]}
使用莫尔斯--威腾复形，
这里使用奇异链复形。
这一差别虽颇具技术性，
但在本书出现多种“高阶”项的情形，
似乎更难像
\textbf{[Oh96I]}
那样使用莫尔斯--威腾复形。
其二，
\textbf{[Oh96I]} 用莫尔斯指标定义滤过，
而这里用能量；
对于单调拉格朗日子流形，
按单调性的定义，
这两个滤过一致。

为了构造谱序列并研究其收敛性，
需要研究万有诺维科夫环上模的一般性质。
由于该环不是诺特环，
构造时必须格外谨慎。

\subsection{谱序列}
\label{subsec:6-1-1}

为陈述本章结果，
简要回顾万有诺维科夫环
$\Lambda_{\mathrm{nov}}(R)$
及其子环
$\Lambda_{0,\mathrm{nov}}(R)$
的基本性质，其中 $R$ 为交换环。
当 $R=\mathbb Q$ 时省去记号中的 $R$，
即写
$\Lambda_{\mathrm{nov}}
=\Lambda_{\mathrm{nov}}(\mathbb Q)$、
$\Lambda_{0,\mathrm{nov}}
=\Lambda_{0,\mathrm{nov}}(\mathbb Q)$。
规定
$\deg(T^\lambda e^n)=2n$，
它们便成为分次环。
分别用
$\Lambda_{\mathrm{nov}}^{(k)}$、
$\Lambda_{0,\mathrm{nov}}^{(k)}$
表示其 $k$ 次部分。
我们已在
$\Lambda_{\mathrm{nov}}$
上定义滤过
\begin{equation}
F^\lambda\Lambda_{\mathrm{nov}}
=
\left\{
\sum_i a_iT^{\lambda_i}e^{n_i}
\ \middle|\ 
\lambda_i\ge\lambda
\right\}.
\tag{6.1.1}\label{eq:6-1-1}
\end{equation}
它诱导
$\Lambda_{0,\mathrm{nov}}$
上的滤过；
这些滤过给两环以完备的一致结构。

上述滤过又在
$\Lambda_{\mathrm{nov}}$
和 $\Lambda_{0,\mathrm{nov}}$
上诱导下列滤过：
\begin{equation}
\mathfrak F^q\Lambda_{\mathrm{nov}}
=F^{q\lambda_0}\Lambda_{\mathrm{nov}},
\qquad
\mathfrak F^q\Lambda_{0,\mathrm{nov}}
=F^{q\lambda_0}\Lambda_{0,\mathrm{nov}}.
\tag{6.1.2}\label{eq:6-1-2}
\end{equation}
其中 $\lambda_0>0$ 稍后确定；
对 $\Lambda_{0,\mathrm{nov}}$
取 $q\in\mathbb Z_{\ge0}$，
对 $\Lambda_{\mathrm{nov}}$
则取 $q\in\mathbb Z$。
$\lambda_0$ 与有隙条件
6.3.16 有关。
相应的伴随分次模为
\begin{equation}
\operatorname{gr}_q
(\mathfrak F\Lambda_{\mathrm{nov}})
=
\frac{\mathfrak F^q\Lambda_{\mathrm{nov}}}
{\mathfrak F^{q+1}\Lambda_{\mathrm{nov}}},
\qquad
\operatorname{gr}_q
(\mathfrak F\Lambda_{0,\mathrm{nov}})
=
\frac{\mathfrak F^q\Lambda_{0,\mathrm{nov}}}
{\mathfrak F^{q+1}\Lambda_{0,\mathrm{nov}}}.
\tag{6.1.3}\label{eq:6-1-3}
\end{equation}

本章的主要结果是导论中的定理 D：
它给出的谱序列\liftIndex{I050}\liftIndex{I005}
描述（奇异）上同调与
弗洛尔上同调之间的关系。
我们将在第 6.4 节证明它。
第 6.3 节先在纯代数语境下
构造谱序列并证明其收敛；
第 6.4.1 小节再于几何情形
验证所需性质，
从而完成 \textup{(D.1)}、
\textup{(D.2)} 的证明。
证明 \textup{(D.3)}
将使用第 3.8 节的算子
$\mathfrak p,\mathfrak q$，
安排在第 6.4.2 小节；
最后于第 6.4.5 小节
证明 \textup{(D.4)}，
完成定理 D。
注意 \textup{(D.3)} 中
$\mathbf b_1=\mathbf b_0$
的附加假设不能去掉
（见注 6.4.12）。

对 $M$ 中一对相对 spin
的拉格朗日子流形
$(L^{(0)},L^{(1)})$，
还可构造博特--莫尔斯版本的谱序列。
下述定理的记号见
引理--定义 3.8.75。
先设这对拉格朗日子流形干净相交，
并写
\[
L^{(1)}\cap L^{(0)}
=
\bigcup_{h\in\pi_0(L^{(1)}\cap L^{(0)})}
R_h.
\]
记 $\mu_L([h,w])$
为马斯洛夫--莫尔斯指标
（见定义 3.7.62）。
$\Theta^-_h$
是由表示
$\pi_1(R_h)\to\{\pm1\}$
诱导的局部系统
（见定义 8.8.2）。
$\pi_{\mathrm{amb}}$ 的记号见
(4.3.27)，
$\mathfrak{PO}$ 见引理 4.3.28；
其余记号见第 3.7.5 小节
的定义 3.7.52 与命题 3.7.59。

\begin{theorem}{6.1.4}\label{thm:6-1-4}
设 $(L^{(0)},L^{(1)})$
为 $M$ 中干净相交的\translatornote{原书实体页 365
在定理句中未重述“干净相交”，
但紧前已以此假设定义各 $R_h$，
且 (6.1.5.1) 必须用到它。
见疑误 A-0353。}
相对 spin 拉格朗日子流形对。
若各
$b_i\in\mathcal M_{\mathrm{weak,def}}(L^{(i)})$
满足
\[
\mathfrak{PO}_{L^{(0)}}(b_0)
=\mathfrak{PO}_{L^{(1)}}(b_1),
\qquad
\pi_{\mathrm{amb}}(b_0)
=\pi_{\mathrm{amb}}(b_1),
\]
则存在具有下列性质的谱序列
$\{E_r^{p,q},\delta_r\}$。
其初页满足
\begin{equation}
E_2^{p,q}
=
\bigoplus_{\substack{
k,\,[h,w]\in
\widetilde{\pi}_0(L^{(1)}\cap L^{(0)})
}}
H^{k-\mu([h,w])}(R_h;\Theta^-_h)
\otimes
\operatorname{gr}_q
(\mathfrak F\Lambda_{0,\mathrm{nov}})^{(p-k)}.
\tag{6.1.5.1}\label{eq:6-1-5-1}
\end{equation}

弗洛尔上同调
$HF((L^{(1)},b_1),(L^{(0)},b_0);
\Lambda_{0,\mathrm{nov}})$
上存在滤过 $\mathfrak F^*$，
使\translatornote{原书实体页 365
本式分母漏写上同调次数 $p$；
导论定理 D 的 \textup{(D.2)}
在同一商式中明确保留 $HF^p$。
见疑误 A-0354。}
\begin{equation}
E_\infty^{p,q}
\cong
\frac{
\mathfrak F^q HF^p
((L^{(1)},b_1),(L^{(0)},b_0);
\Lambda_{0,\mathrm{nov}})
}{
\mathfrak F^{q+1} HF^p
((L^{(1)},b_1),(L^{(0)},b_0);
\Lambda_{0,\mathrm{nov}})
}.
\tag{6.1.5.2}\label{eq:6-1-5-2}
\end{equation}

\par\smallskip\noindent\textup{(6.1.5.3)}\quad
这个谱序列与
$HF((L^{(1)},b_1),(L^{(0)},b_0);
\Lambda_{0,\mathrm{nov}})$
在自交弗洛尔上同调对
\[
HF((L^{(j)},b_j),(L^{(j)},b_j);
\Lambda_{0,\mathrm{nov}}),
\qquad j=0,1,
\]
上的双模结构相容；
该结构定义于第 3.8.10 小节。
具体说，滤过 $\mathfrak F$
与双模结构相容；
本定理的 $E_r$
是定理 D 中相应 $E_r$ 上的双模，
且 $\delta_r$ 是导子；
\textup{(6.1.5.1)}、
\textup{(6.1.5.2)}
中的同构都是双模同构。
\end{theorem}

为证明定理 D 和 6.1.4，
将在第 6.3 节建立代数结果，
再分别应用于与一个拉格朗日子流形
及一对拉格朗日子流形
相联系的复形\translatornote{原书实体页 365
在第二个复形中把双模微分下标写成
$(b_0,b_1)$；
本书双模左右输入及本页
(6.1.5.2)--(6.1.5.3)
均按 $(b_1,b_0)$ 排列。
同处 $\Lambda_{0.nov}$ 的句点
也按全书记号写成逗号。
见疑误 A-0355。}
\[
(C(L;\Lambda_{0,\mathrm{nov}}),
\mathfrak m^{\mathbf b,b}_1)
\quad\text{和}\quad
(C(L^{(1)},L^{(0)};
\Lambda_{0,\mathrm{nov}}),
\delta_{b_1,b_0}).
\]

\begin{remark}{6.1.6}\label{rem:6-1-6}
$\Lambda_{\mathrm{nov}}$
是 $\Lambda_{0,\mathrm{nov}}$
的局部化，因此在后者上平坦。
还应注意，谱序列的
$E_\infty$ 项将用逆极限
来定义（定义 6.3.23）。
不过，定理 6.3.28 将表明：
若复形是弱有限的微分分次模，
则存在 $r_0$，
使
\[
E_{r_0}^{p,q}
\cong E_{r_0+1}^{p,q}
\cong\cdots
\cong E_\infty^{p,q}.
\]
因此，也可通过把系数环改为
$\Lambda_{\mathrm{nov}}$，
在该环上构造定理 D
及 6.1.4 的谱序列。
还记得我们在第 5.3 节
已证明（见定理 4.1.5）
\begin{equation}
\begin{aligned}
&HF((L,(\mathbf b,b)),\psi;
\Lambda_{\mathrm{nov}})\\
&\quad\cong
HF((L,(\mathbf b,b)),
(L,(\mathbf b,b));
\Lambda_{0,\mathrm{nov}})
\otimes_{\Lambda_{0,\mathrm{nov}}}
\Lambda_{\mathrm{nov}},
\end{aligned}
\tag{6.1.7}\label{eq:6-1-7}
\end{equation}
其中 $\psi$ 为哈密顿微分同胚，且
\begin{equation}
\begin{aligned}
&HF((L,(\mathbf b,b)),\psi;
\Lambda_{\mathrm{nov}})\\
&\quad=
HF\bigl(
(L,(\mathbf b,b)),
(\psi(L),(\psi_*\mathbf b,\psi_*b));
\Lambda_{\mathrm{nov}}
\bigr).
\end{aligned}
\tag{6.1.8}\label{eq:6-1-8}
\end{equation}
这是 $\Lambda_{\mathrm{nov}}$
系数的弗洛尔上同调。
(6.1.7) 与 \textup{(D.3)}
蕴含导论中的定理 H
（见第 6.5.1 小节）。
\end{remark}
